% Part: proof-theory
% Chapter: cut-elimination
% Section: introduction

\documentclass[../../../include/open-logic-section]{subfiles}

\begin{document}

\olfileid{pt}{cut}{int}

\olsection{પરિચય}

\Cut{} નિયમ આ છે:
\[
\Axiom$\Gamma \fCenter \Delta, !A$
\Axiom$!A, \Pi \fCenter \Lambda$
\RightLabel{\Cut}
\BinaryInf$\Gamma, \Pi \fCenter \Delta, \Lambda$
\DisplayProof
\]
(!!{formula}~$!A$ને
\emph{કટ-સૂત્ર} કહે છે.)

\Cut{} એ ગેન્ઝેનના મૂળ
સિક્વન્ટ કલન~\Log{LK}માં હતું.
ગેન્ઝેનનું \emph{Hauptsatz}
(``મુખ્ય પ્રમેય'') કહે છે કે
~\Log{LK}માં !!{provable}
દરેક સિક્વન્ટ \Cut{} નિયમ વિના
પણ !!{provable} છે;
એટલે~\Log{LK}-સાબિતીઓમાંથી
\Cut{} ``દૂર'' કરી શકાય.
આપણી પ્રણાલીઓ~\Log{G1c}
અને~\Log{G3c}માં \Cut{} નથી,
પણ તેમના માટે પણ પ્રશ્ન ઊભો થાય:
\Log{G1c}ના (અથવા~\Log{G3c}ના)
બીજા નિયમો સાથે~\Cut{} વાપરતી
!!{proof}sમાંથી \Cut{} દૂર કરી શકાય?
પહેલેથી નક્કી કરેલી પરિભાષામાં
એ જ પ્રશ્ન છે: શું~\Cut{}
એ~\Log{G1c}નો
(અથવા~\Log{G3c}નો)
સ્વીકાર્ય નિયમ છે?

કટ-નિરાકરણ પ્રમેય સાબિતી સિદ્ધાંતનું
સંભવતઃ સૌથી મહત્ત્વનું અને જાણીતું
પરિણામ છે. પહેલાં તો તે બતાવે છે કે
~\Log{G1c} અને~\Log{G3c}ને
પ્રથમ નજરે અનિવાર્ય લાગે એવો નિયમ
જરૂરી નથી. વાસ્તવિક સાબિતીઓનું
ઔપચારિકીકરણ કરો ત્યારે ઉપપ્રમેયો
વાપરવાની રીત જોઈએ જ.
ગણિતજ્ઞો એક પરિણામ સાબિત કરે
અને પછી બીજી સાબિતીઓમાં તેનો
આધાર લે છે. તેથી પ્રથમ પરિણામ~$L$
(ઉપપ્રમેય)ને
$\Sequent L$ સિક્વન્ટની
!!{proof} તરીકે ઔપચારિક કરી શકાય.
પછી~$L$ વાપરતા બીજા
પરિણામ~$R$ની સાબિતીને
$L \Sequent R$ સિક્વન્ટની
!!{proof} તરીકે લખો.
માત્ર~$R$ની સાબિતી,
એટલે $\Sequent R$ની
!!a{proof}, છે તે કેવી રીતે જાણવું?
બે !!{proof}s જોડવાની રીત
જોઈએ; \Cut{} નિયમ એ કરવા દે છે.

આપણે કહ્યું હતું કે~\Log{G1c}
અને~\Log{G3c} પૂર્ણ છે;
એટલે સાબિત થવા જોઈએ એવા
બધા પ્રામાણ્ય ધરાવતા સિક્વન્ટ
તેઓ સાબિત કરે છે. આ સીધું
પણ સાબિત કરી શકાય.
પરંતુ ગેન્ઝેન લખતા હતા ત્યારે
માત્ર સ્વયંસિદ્ધિમૂલક નિષ્પત્તિ
પ્રણાલીઓની પૂર્ણતા જાણીતી હતી.
~\Log{LK}ની પૂર્ણતા બતાવવા
તેમણે સ્વયંસિદ્ધિમૂલક પ્રણાલીની
!!{derivation}sને
\Log{LK}-!!{proof}sમાં રૂપાંતરિત
કરવાની રીત આપી. તેમાં~\Cut{}
અનિવાર્ય રીતે વપરાતો હતો.
કટ-નિરાકરણ માત્ર શાસ્ત્રીય તર્ક
માટે મહત્ત્વનું નથી: બીજા અનેક
તર્કો માટે પણ સિક્વન્ટ કલન છે.
કેટલાક તર્કોના સિક્વન્ટ કલનની
પૂર્ણતા સીધી સાબિત કરવી મુશ્કેલ
અથવા અશક્ય હોય છે; તેથી
બીજી પ્રણાલીઓમાંથી~\Cut{}
વાપરતા રૂપાંતરો જ પૂર્ણતા
બતાવવાનો માર્ગ હોય છે.
પછી~\Cut{} બિનજરૂરી છે અને
તે વિનાની પ્રણાલી પૂર્ણ છે
એ સ્થાપવા સાબિતી-સૈદ્ધાંતિક
રીતે કટ-નિરાકરણ સાબિત કરવું પડે.

કટ-નિરાકરણ વડે સ્વતંત્ર રસનાં
બીજાં પરિણામો પણ મેળવી શકાય
એથી પણ તે મહત્ત્વનું છે.
આપણે ક્રેગના અંતર્વેશન ઉપપ્રમેય
અને હર્બ્રાન્ડના પ્રમેયને જોઈશું.

કટ-નિરાકરણની સાબિતીમાં સામાન્ય રીતે
~\Cut{} વાપરતી !!a{proof}ને
તે વિનાની સાબિતીમાં કેવી રીતે
ફેરવવી તે બતાવાય છે.
આ રૂપાંતરો મોટા ભાગે
``સ્થાનિક'' હોય છે:
!!a{proof}નો એક ભાગ
(સામાન્ય રીતે~\Cut{} અનુમાનમાં
પૂરી થતી ઉપ-!!{proof}) લઈને
તેને એ જ સિક્વન્ટની બીજી
ઉપ-!!{proof}થી બદલીયે,
જેમાં~\Cut{} અનુમાન દૂર થયું
હોય અથવા બીજા અનુમાનોથી
બદલાયું હોય.
\sourcecorrection{OLMD-213}{મૂળમાં glossary ચિહ્ન !!a{proof} બે પંક્તિ વચ્ચે તૂટેલું છે; અહીં અખંડ ચિહ્ન પુનઃસ્થાપિત કર્યું છે.}
આવી દલીલો
\Cut{} ધરાવતી !!{proof}sને
તે વિનાની સાબિતીઓમાં ફેરવવાના
પગલું-દર-પગલું કલનવિધિ આપે છે.
ઉપયોગોમાં કટ-નિરાકરણ પ્રક્રિયાઓ
લઈએ ત્યારે આ વધુ માહિતી આપે છે.
દાખલા તરીકે, અંતર્વેશન ઉપપ્રમેયની
નિદર્શ-સૈદ્ધાંતિક સાબિતી કહી શકે
કે જો~$!A \lif !B$ પ્રામાણ્ય ધરાવે
તો અંતર્વેશક છે
(હમણાં અંતર્વેશક શું છે તે છોડીએ).
કટ-નિરાકરણ વાપરતી
સાબિતી-સૈદ્ધાંતિક દલીલ
$!A \lif !B$ની !!a{proof}
લઈ અંતર્વેશક આપતો કલનવિધિ આપે છે.
નિદર્શ-સૈદ્ધાંતિક દલીલની તુલનામાં
આ દલીલ \emph{રચનાત્મક} છે.

કટ-નિરાકરણ પ્રમેય સાબિત કરવાની
બે પ્રચલિત રીતો છે.
પહેલી (ગેન્ઝેનથી શરૂ થયેલી)
સાબિતીમાં સૌથી ઉપરના કટ દૂર કરી
શકાય છે એ બતાવે છે, અને
પછી એક પછી એક એવા કટ
દૂર કરે છે. બીજી
(મૂળે શ્યુટે અને ટાઇટની)
!!a{proof}માં સૌથી જટિલ કટ
પર ધ્યાન આપે છે અને એવા
મહત્તમ કટ ક્રમશઃ દૂર કરે છે
જેનાથી વધુ ઊંડાઈનું કટ
!!{formula} બીજા કટમાં ન હોય.
દરેક રીતના લાભ અને મર્યાદાઓ છે.
કેટલીક પ્રણાલીઓમાં એક રીત ચાલે,
જ્યારે બીજી અમલમાં મૂકવી
મુશ્કેલ કે અશક્ય હોય.
તેથી બંને રીતો વર્ણવીશું.
કટ-નિરાકરણ~\Log{G3c} માટે
સાબિત કરીશું. વાસ્તવમાં
~\Cut{}ને નહીં, પણ
\emph{સંદર્ભ-સહભાગી કટ}
~\CutCS{}ને દૂર કરી શકાય
એ બતાવીશું:
\[
\Axiom$\Gamma \fCenter \Delta, !A$
\Axiom$!A, \Gamma \fCenter \Delta$
\RightLabel{\CutCS}
\BinaryInf$\Gamma \fCenter \Delta$
\DisplayProof
\]
\sourcecorrection{OLMD-214}{મૂળ બીજા વૃક્ષનું વર્ણન $\CutCS$નું છે, પણ વૃક્ષનું નિયમ-લેબલ $\Cut$ છે; સમાન સંદર્ભવાળા કટનું યોગ્ય લેબલ અહીં પુનઃસ્થાપિત કર્યું છે.}
જો સિક્વન્ટ કલન શિથિલીકરણ અને
સંકોચનને મંજૂરી આપે
(~\Log{G1c}ની જેમ સ્પષ્ટ
નિયમો તરીકે કે~\Log{G3c}ની
જેમ સ્વીકાર્ય નિયમો તરીકે),
તો~\Cut{} વડે~\CutCS{}નું
અને તેનાથી વિરુદ્ધ પણ
અનુકરણ કરી શકાય.
અમે~\Cut{}ને બદલે~\CutCS{}
પસંદ કરીએ છીએ, કારણ કે પહેલી
કટ-નિરાકરણ રીત~\CutCS{} માટે
વર્ણવવી સહેલી છે અને બીજી
રીત માત્ર~\CutCS{} માટે ચાલે છે.

\end{document}
